

\chapter{Related Work}
\label{chap:related_work}

Electricity demand forecasting has received substantial attention in the literature,
yet studies addressing large-scale distributed commercial networks, particularly those
combining differing site characteristics, EV charging infrastructure, and real market
penalty structures remain scarce.

Research on electricity demand forecasting using Ukrainian data is notably limited in
both volume and methodological scope. The most directly relevant study is by
\citet{sinitsyn2023}, conducted at the Institute of Software Systems of the National
Academy of Sciences of Ukraine, which addresses forecasting of electricity generation
and consumption using machine learning. The authors applied seven regression algorithms
from the scikit-learn library --- including k-nearest neighbors (KNN), Random Forest, Gradient
Boosting, and linear regression --- to hourly data from Ukraine's four electricity market
segments: bilateral contracts, day-ahead, intraday, and balancing markets. For the task
of reconstructing missing bilateral contract data, the best-performing KNN configuration
achieved $R^2 = 0.987$ and Mean Absolute Percentage Error (MAPE) of approximately $0.85\%$, leading the authors to
conclude that one-hour-ahead market volume forecasting can reach accuracy near $1\%$.
While this represents a useful proof of concept for data-driven approaches on Ukrainian
market data, the study has several fundamental limitations for the purposes of distributed
site-level forecasting. The forecasting target is aggregate market volume --- a
substantially easier task than predicting consumption at hundreds of individual
heterogeneous locations. No deep learning architectures were evaluated, the forecasting
horizon was limited to a single step ahead, and the evaluation relied exclusively on
statistical metrics without any assessment of economic value or imbalance cost
implications.

The international literature provides a much richer methodological landscape, beginning
with classical statistical approaches that remain important reference points.
\citet{cai2019} provided a thorough comparison of SARIMAX against deep learning
alternatives for day-ahead building-level load forecasting across three commercial
buildings with different characteristics and load scales (40--650 kW peak). Their results
demonstrated that SARIMAX performs adequately when weather covariance is weak, but
performance deteriorates substantially when temperature rises and air conditioning loads
increase. This highlights traditional models' limitations: they struggle with
nonlinearity and do not naturally scale to large sets of heterogeneous series.

The main disadvantage of ARIMA with exogenous variables and SARIMAX is the need to train a separate model for
each time series. This concept is most directly and explicitly articulated by
\citet{grabner2023}, who state that traditional load forecasting techniques
``commonly learn one prediction model for each given time series and are therefore
considered to be local,'' and explicitly identify that such local models are ``less
suitable for forecasting larger groups of time series, as is the case with distribution
networks.''

Machine learning methods such as Support Vector Regression (SVR), Random Forest and XGBoost
partially overcome these limitations and have shown reasonable performance in electricity
demand forecasting~\cite{sinitsyn2023}. However, these models still often require
substantial feature engineering and are typically applied either to single-series problems
or to pooled tabular settings without explicitly modeling temporal structure across many
locations~\cite{liu2025}. \citet{kyriakopoulos2025}, who systematically evaluated ARIMA,
XGBoost, Gated Recurrent Unit (GRU), Long Short-Term Memory (LSTM), and Transformer models across four EV charging datasets, showed that
ARIMA and XGBoost were competitive only in localized scenarios and did not scale reliably
to longer horizons or higher aggregation levels.

Considering this, \citet{cai2019} proposed a Gated Convolutional Neural Network ((GCNN24)) model for day-ahead building-level
forecasting that outperformed both SARIMAX and gated RNN alternatives, improving accuracy
by $22.6\%$ over SARIMAX while achieving a $25$--$42\%$ speedup in computation time
compared to Recurrent Neural Network (RNN) models. The GCNN24 model proved superior to
recursive one-hour models because it learns daily load shapes holistically rather than
accumulating single-step prediction errors.

The most significant benefit of deep learning approaches compared to classical models is
the move from local models (one model per time series) to global models (one model for
many time series). \citet{grabner2023} articulate the fundamental limitations of local
modeling for distribution networks: local models ignore dependencies between loads,
require sufficient historical data for each individual consumer, cannot handle locations
with missing history, and are incapable of adapting to changing load behavior without
retraining. Their proposed global framework, built on the Neural Basis Expansion Analysis for Time Series (N-BEATS) deep learning
architecture, demonstrated that a single global model outperformed naive benchmarks by
more than $25\%$ in Mean Absolute Error (MAE) across consumer aggregations ranging from individual households
to large transformer stations, while requiring approximately $30\times$ fewer parameters
than the equivalent set of local models. Training was $3.3\times$ faster for 1,000 time
series compared to training local models individually. Global models can additionally
generate reasonable forecasts for time series with no historical data --- a capability
impossible with local approaches and demonstrated robust performance when load
behavior changed over time, because the model had learned diverse patterns from the
broader training set.

The most widely used architecture in related research is the LSTM network. In the context of residential demand forecasting, \citet{bayer2024} compare standalone LSTM models with CNN–LSTM hybrids, where a Convolutional Neural Network is employed for feature extraction and the LSTM component captures temporal dependencies. Their study benchmarks these approaches against synthesized load profiles (SLPs), which remain standard practice among many small and medium-sized utilities. Using smart meter data from 3,511 households in Germany, the authors demonstrate that LSTM-based approaches outperform SLPs by up to $68.5%$ in terms of Root Mean Square Error (RMSE) for day-ahead forecasts. Importantly, they also evaluated performance under simulated future grid states with increased distributed photovoltaic (PV) generation, finding that all methods performed worse under
higher PV penetration.

Transformer architectures represent the most recent frontier in load forecasting. Their
self-attention mechanism enables direct modeling of relationships between any positions
in a sequence without the sequential processing bottleneck of recurrent networks.
\citet{moveh2025a} developed a modified Transformer framework with hierarchical attention
mechanisms for multi-building energy prediction across 100 commercial buildings in three
climate zones. Their architecture processes building-specific features and cross-building
correlations through separate attention levels before combining them through cross-attention
for final prediction. The framework achieved MAPE of $3.2\%$, a $23.7\%$ improvement
over traditional methods with stable performance across building types (offices:
$2.8\%$, healthcare: $3.5\%$) and near-linear computational scaling with building count.
The hierarchical attention design is particularly instructive: it allows the model to
learn both what is unique to each location and what is shared across the network, which
is precisely the modeling challenge in a multi-site commercial forecasting system.

\section{Evaluation and Loss Functions}

The reviewed literature overwhelmingly relies on MAE, RMSE, and MAPE. These metrics are
useful for comparison, but they have important limitations. \citet{koponen2020} argue that
conventional error metrics often fail to reflect the real economic consequences of forecast
errors, particularly because under-forecasting and over-forecasting may have asymmetric
costs, and because peak-load errors are often more expensive than off-peak errors. Moreover,
many of these metrics have significant flaws: MAPE penalizes over-forecasts more than
under-forecasts, gives relatively little weight to large absolute errors during peak loads,
and is undefined for zero values; MAE assumes linear cost dependence on error magnitude;
and RMSE may or may not reflect actual cost dynamics. Through simulation of Finnish
balancing market costs, they showed that forecasting error costs are highly stochastic,
dominated by rare high-price peaks, and quadratically rather than linearly related to
error magnitude when errors are small relative to the total balance.

\citet{zhang2021} proposed a generalized cost-oriented load forecasting framework that
directly addresses this gap, showing that models trained or evaluated with economically
meaningful asymmetric loss functions can outperform models with conventional error metrics
in real operational value. An Artificial Neural Network (ANN) model with an hourly cost-oriented loss achieved up to
$13.74\%$ improvement in economic costs compared to the Mean Square Error (MSE) trained benchmark, and a Multiple Linear Regression
(MLR) model with the same cost-oriented loss improved economic costs by $10.19\%$. Both
cost-oriented models achieved better economic outcomes despite slightly higher statistical
forecasting errors.

\section{EV Charging: A Special Case}
\label{sec:ev_charging}

The growing presence of EV chargers at commercial locations introduces a forecasting
challenge that is qualitatively different from conventional load prediction. EV charging
demand is driven by individual user decisions arrival time, battery state, dwell time
 rather than by physical processes like Heating, Ventilation, and Air Conditioning cycling or lighting schedules, making it
inherently more stochastic, intermittent, and spiky. Charging station time series contain
numerous zero-valued intervals interspersed with sharp demand peaks, a pattern that
conventional forecasting methods struggle to model effectively~\cite{tang2024,
koohfar2023}.

\citet{vanetten2024} proposed a framework for large-scale EV charging demand forecasting
across multiple stations using global deep learning. Their Neural Hierarchical Interpolation for Time Series (N-HiTS) model trained globally
on data from multiple stations consistently outperformed locally trained counterparts,
ARIMA, naive baselines, and even Transformer models across four datasets from London,
Palo Alto, Perth, and Boulder. The global N-HiTS model reduced MAE by $51.6\%$ compared
to locally trained N-HiTS on London data. The paper highlights a practical evaluation
concern: high MAPE values resulted from the low signal-to-noise ratio in individual
station data, where near-zero consumption periods inflate percentage errors. This means
that error rates for EV charger data are quite high, with the lowest MAPE for day-ahead
forecasting at $30.16\%$ and for 30 days at $42.41\%$, highlighting that accurately
forecasting EV charging loads is substantially harder than predicting conventional
consumption.

\section{Feature Engineering and Heterogeneity}
\label{sec:feature_engineering_literature}

Across the reviewed literature, certain input features consistently emerge as important.
Temperature is the dominant weather variable, showing correlations of $0.51$--$0.76$
with building load across different commercial buildings~\cite{cai2019}.
\citet{tian2025} performed more granular analysis using both Pearson correlation and
Maximal Information Coefficient, confirming that temperature and humidity exhibit the
strongest relationships with load, with the strength varying seasonally. Calendar
features --- hour, day of week, month, holidays --- appear universally across
studies~\cite{sinitsyn2023, ye2024, moveh2025a}.

Handling heterogeneity across locations remains a practical challenge. \citet{spencer2024}
demonstrated that dataset characteristics critically influence model performance, with
weather feature consistency being more important than geographic or climate zone
similarity. \citet{grabner2023} addressed heterogeneity through clustering and
localization, while \citet{moveh2025a} used building-type-specific normalization
parameters. For weather forecast uncertainty specifically, \citet{xu2025} showed that
one-day-ahead probabilistic weather forecasts produced the most accurate probabilistic
load predictions compared to longer-range forecasts, and \citet{cai2019} demonstrated
that their GCNN24 model's RMSE increased by only $3.36\%$ even when $60\%$ of weather
forecast data was corrupted with noise, indicating meaningful robustness to weather
prediction errors.

\section{Model Benchmarks}
\label{sec:benchmarks}

To situate candidate architectures within the broader landscape, two large-scale
benchmarking studies were reviewed. \citet{qiu2024tfb} evaluated a wide range of models
--- from ARIMA and Linear Regression through RNNs to several Transformer variants --- on
multivariate time series including the Electricity dataset, containing hourly electricity
consumption in kWh from 2012 to 2014. PatchTST displayed the best performance, with
Crossformer not far behind. \citet{li2024tsfmbench} focused on foundation models for time
series forecasting, with Chronos and ROSE showing the best short-term and mid-term
forecasting performance on the Electricity dataset, and MOIRAI displaying the best
long-term performance.

\section{Summary}

Overall, while classical and machine learning approaches can be useful in certain cases,
multiple studies outline their limitations: limited scalability for ARIMA-family models,
and a lack of ability to model temporal dynamics~\cite{kyriakopoulos2025}. For consumption
prediction, deep learning models and especially LSTM-based architectures are most
prominent. Newer Transformer-based models featuring attention mechanisms have seen only
limited usage and are not always superior to strong deep learning alternatives, with
\citet{vanetten2024} showing that N-HiTS can outperform Transformer models. Regarding
loss functions, \citet{zhang2021} showed that using a loss function reflecting the real
economic costs of prediction errors can yield up to $13\%$ in savings compared to an
equivalent model trained for statistical accuracy. Another important finding is that EV
load forecasting carries significantly higher error rates than conventional load
forecasting, due to the low signal-to-noise ratio~\cite{vanetten2024}. Feature engineering
is also a critical component of model training, with weather and calendar features
appearing across many studies~\cite{sinitsyn2023, ye2024, moveh2025a}. Notably,
\citet{cai2019} showed that even when $60\%$ of weather data was corrupted by noise, RMSE
increased by only $3.36\%$, highlighting meaningful model robustness to weather forecast
errors. Finally, a review of relevant benchmarks identified the best-performing
Transformer models on the Electricity dataset as PatchTST, Crossformer, Chronos, ROSE,
and MOIRAI~\cite{li2024tsfmbench, qiu2024tfb}.